\chapter[Computational Demands of Visual Reasoning]{Understanding the computational demand underlying visual reasoning}
\label{chapter1a}

\startcontents[chapters]
\printmyminitoc

\section{Introduction}

Humans can effortlessly reason about the visual world and provide rich and detailed descriptions of briefly presented real-life photographs \citep{Fei-Fei2007-zb}, vastly outperforming the best current computer vision systems \citep{Geman2015-jm, Kreiman2020-zd}. 
For the most part, studies of visual reasoning in humans have sought to characterize the neural computations underlying the judgment of individual relations between objects, such as their spatial relations (e.g., \citet{logan1994ability}) or whether they are the same or different (up to a transformation, e.g., \citet{shepard1971mental}). It has also been shown that different visual reasoning problems have different attentional and working memory demands \citep{logan1994spatial,Moore1994,Rosielle2002,Holcombe2011,vanderham2012,kroger2002recruitment,Golde2010,clevenger2014working,brady2015contextual}. However, there is still little known about the neural computations that are engaged by different types of visual reasoning (see~\citet{ricci37same} for a recent review). 

One benchmark that has been designed to probe abstract visual relational capabilities in humans and machines is the \textit{Synthetic Visual Reasoning Test} (SVRT) ~\citep{fleuret2011comparing}. The dataset consists of twenty-three hand-designed binary classification problems that test abstract relationships between objects posed on images of closed-contour shapes. Observers are never explicitly given the underlying rule for solving any given problem. Instead, they learn it while classifying positive and negative examples and receiving task feedback. Examples from two representative tasks are depicted in Figure~\ref{fig:example}: observers must learn to recognize whether two shapes are the same or different (Task \textit{1}) or whether or not the smaller of the two shapes are near the boundary (Task \textit{2}). Additional abstract relationships tested in the challenge include ``inside", ``in between”, ``forming a square”, ``aligned in a row" or ``finding symmetry" (see Figures \ref{fig:exampleSD} and \ref{fig:exampleSR} for examples).

\begin{figure}[t]
\centering
  \includegraphics[width=1\linewidth]{images/chapter1/svrtexample3}
 \caption{Two SVRT sample tasks from a set of twenty-three in total. For each task, the leftmost and rightmost two examples illustrate the two categories to be classified. Representative samples for the complete set of twenty-three tasks can be found in Figure~\ref{fig:exampleSD} and~\ref{fig:exampleSR}.} \label{fig:example}
\end{figure}


Most SVRT tasks are rapidly learned by human observers within twenty or fewer training examples~\citep{fleuret2011comparing} (see Table~\ref{table:humans}; reproduced from the original study). On the other hand, modern deep neural network models require several orders of magnitude more training samples for some of the more challenging tasks~\citep{ellis2015unsupervised,kim2018not,MESSINA2021,stabinger2021evaluating,Stabinger2016,Puebla2021.04.06.438551} (see~\citet{ricci37same} for review; see also \citet{Funke2021} for an alternative perspective). 

It is now clear that some SVRT tasks are more difficult to learn than others. For instance, tasks that involve spatial-relation (SR) judgments can be learned much more easily by deep convolutional neural networks (CNNs) than tasks that involve same-different (SD) judgments~\citep{Stabinger2016, kim2018not,yihe2019program}. In contrast, a very recent study \citep{Puebla2021.04.06.438551} demonstrated that even when CNNs learn to detect whether objects are the same or different, they fail to generalize over small changes in appearance, meaning that they have only partially learned this abstract rule. The implication of the relative difficulty of learning SR versus SD tasks is that CNNs appear to need additional computations to solve SD tasks beyond standard filtering, non-linear rectification, and pooling. Indeed, recent human electrophysiology work~\citep{AlamiaENEURO.0267-20.2020} has shown that SD tasks recruit cortical mechanisms associated with attention and working memory processes to a greater extent than SR tasks. \textcolor{black}{Others have argued that SD tasks are central to human intelligence \citep{firestone2020performance, FORBUS202163, GENTNER202184}.} Beyond this basic dichotomy of SR and SD tasks, little is known about the neural computations necessary to learn to solve SVRT tasks as efficiently as human observers.


Here, we investigate the neural computations required for visual reasoning. In our experiment, we extend prior studies on the learnability of individual SVRT tasks by feedforward neural networks using a popular class of deep neural networks known as deep residual networks (``ResNets'') \citep{he2016deep}. We systematically analyze the ability of ResNets to learn all twenty-three SVRT tasks as a function of their expressiveness, parameterized by processing depth (number of layers), and their efficiency in learning a particular task. Through these experiments, we found that most of the performance variance in the space of SVRT tasks could be accounted for by two principal components, which reflected both the type of task (same-different vs. spatial-relation judgments) and the number of relations used to compose the underlying rules. 

\section{Systematic analysis of SVRT tasks' learnability}


All experiments were carried out with the \textit{Synthetic Visual Reasoning Test} (SVRT) dataset using code provided by the authors to generate images with dimension \textit{128 $\times$ 128} pixels (see \citet{fleuret2011comparing} for details). \textcolor{black}{ All images were normalized and resized to 256$\times$256 pixels for training and testing models. No image augmentations were used during training.}
In our first experiment, we wanted to measure how easy or difficult each task is for ResNets to learn. We did this by recording the SVRT performance of multiple ResNets, each with different numbers of layers and trained with different numbers of examples. By varying model complexity and the number of samples provided to a model to learn any given task, we obtained complementary measures of the learnability of every SVRT task for ResNet architectures. In total, we trained 18-, 50-, and 152-layer ResNets separately on each of the SVRT's twenty-three tasks. Each of these models was trained with .5k, 1k, 5k, 10k, 15k, and 120k class-balanced samples. We also generated two unique sets of 40k positive and negative samples for each task: one was used as a validation set to select a stopping criterion for training the networks (if validation accuracy reaches 100\%) and one was used as a test set to report model accuracy. In addition, we used three independent random initializations of the training weights for each configuration of architecture/task and selected the best model using the validation set. Models were trained for \textit{100} epochs using the $Adam$ optimizer \citep{kingma2014adam}  with a training schedule (we used an initial learning rate of 1$e$-3 and changing it to 1$e$-4  from the $70^{th}$ epoch \textcolor{black}{onward}). As a control, because these tasks are quite different from each other, we also tested two additional initial learning rates (\textit{1e-4, 1e-5}). 


Consistent with prior work \citep{kim2018not,Stabinger2016,yihe2019program}, we found that some SVRT tasks are much easier to learn than others for ResNets (Figure~\ref{fig:overall}). For instance, a ResNet50 needs only \textit{500} examples to perform well on tasks \textit{2, 3, 4, 8, 10, 11, 18} but the same network needs \textit{120k} samples to perform well on task \textit{21} (see Figures~\ref{fig:exampleSD} and \ref{fig:exampleSR} for examples of these tasks). Similarly, with \textit{500} training examples, task \textit{2, 3, 4 \& 11} can be learned well with only 18 layers while task \textit{9, 12, 15 \& 23} require as many as 152 layers. A key assumption of our work is that these differences in training set sizes and depth requirements between different SVRT tasks reflect different computational strategies that need to be discovered by the neural networks during training for different tasks. Our next goal is to characterize what these computational strategies are.

\begin{figure}[htbp]
\centering
  \includegraphics[width=.5\textwidth]{images/chapter1/resnet18_2} \includegraphics[width=.5\textwidth]{images/chapter1/resnet50_2} \\
  \includegraphics[width=.5\textwidth]{images/chapter1/resnet152_2}
 \caption{Test accuracy for each of the twenty-three SVRT tasks as a function of the number of training samples for ResNets with depths 18, 50 and 152, resp. The color scheme reflects the identified taxonomy of SVRT tasks (see Figure~\ref{fig:clustering} and text for details).}\label{fig:overall}
\end{figure}


\section{An SVRT taxonomy}
To better understand the computational strategies needed to solve the SVRT, we analyzed ResNet performance on the tasks with a multi-variate clustering analysis. For each individual task, we created an $N$-dimensional vector by concatenating the test accuracy of all ResNet architectures ($N = 3$ depths $\times$ 5 training set sizes = 15), which served as a signature of each task's computational requirements.  We then passed a matrix of these vectors to an agglomerative hierarchical clustering analysis (Figure~\ref{fig:clustering}) using the $Ward's$ method.


\begin{figure}[t!]
\centering
  \includegraphics[width=.95\linewidth]{images/chapter1/dendrogram4}
\caption{Dendrogram derived from an N-dim hierarchical clustering analysis on the test accuracy of N=15 ResNets[18/50/152] trained to solve each task  over a range of training set sizes.} \label{fig:clustering}
\end{figure}



Our clustering analysis revealed a novel taxonomy for the SVRT.
At the coarsest level, it recapitulated the dichotomy between \textit{same-different} (SD; green branches) and \textit{spatial-relation} (SR; \textcolor{black}{brown} branches) categorization tasks originally identified by \citet{kim2018not} using shallow CNNs. Interestingly, two of the tasks which were classified as SR by \citet{kim2018not} (tasks \textit{6 \& 17}) were assigned to the SD cluster in our analysis. We examined the descriptions of these two tasks as given in~\citet{fleuret2011comparing} (see also Figures~\ref{fig:exampleSD} and~\ref{fig:exampleSR}) and found that these two tasks involve both SR and SD: they ask observers to tell whether shapes are the same or different and judge the distance between the shapes. Specifically, task \textit{6} involves two pairs of identical shapes with one category having the same distance in-between two identical shapes vs. not in the other. Similarly, in task \textit{17}, three of the four shapes are identical and their distance from the non-identical one is the same in one category vs. different in the other. Thus, our data-driven dichotomization of SR vs. SD refines the original proposal of \citet{kim2018not}. This could be due to our use of ResNets (as opposed to vanilla CNNs), deeper networks, and a greater variety of training set sizes (including much smaller training set sizes than those used by \citet{kim2018not}). \textcolor{black}{The analysis by \citet{fleuret2011comparing} also revealed that several SD tasks (\textit{6, 16, 17, 21}) are particularly challenging for human observers.}


Our clustering analysis also revealed a finer organization than the main SR vs. SD dichotomy. The SR cluster could be further subdivided into two sub-clusters. The $SR_2$ (dark-brown-coloured) branch in Figure~\ref{fig:clustering} captures tasks that involve relatively simple and basic relation rules such as shapes making close contact (\textit{3, 11}), or being close to one another (\textit{2}), one shape being inside the other (\textit{4}) or whether the shapes are arranged to form a symmetric pattern (\textit{8, 10, 18}). In contrast, tasks that fall in the $SR_1$ (light-brown-colored) branch involve the composition of more than two rules such as comparing the size of multiple shapes to identify a subgroup before identifying the relationship between the members of the sub-groups. This includes tasks such as finding a \textit{larger} object \textit{in between} two smaller ones (\textit{9}) or three shapes of which two are small and one large with two smaller (\textit{identification of large and small object}) ones either inside or outside in one category vs. one \textit{inside} and the other \textit{outside} in the second (\textit{23}), or \textit{two small} shapes \textit{equally close} to a bigger one (\textit{12}), etc. These tasks also tend to be comparatively harder to learn, requiring ResNets with greater processing depth and more training samples. For instance, tasks \textit{9, 12, 15, 23} were harder to learn than tasks \textit{2, 4, 11} requiring more samples and/or more depth to solve well (Figure~\ref{fig:overall}).



We found that task \textit{15} gets assigned to this latter sub-cluster because the task requires finding four shapes in an image that are identical vs. not. One would expect this task to fall in the SD cluster but we speculate that the deep networks are actually able to leverage a shortcut \citep{geirhos2020shortcut} by classifying the overall pattern as symmetric/square (when the four shapes are identical) vs. trapezoid (when the four shapes are different; see Figure~\ref{fig:exampleSR}) -- effectively turning an SD task into an SR task. 


Our clustering analysis also reveals a further subdivision of the SD cluster. These tasks require recognizing shapes that are identical to at least one of the other shapes in the image. The first sub-cluster $SD_2$ (light green color branch) belongs to tasks that require identification of simple task rules, like answering whether or not two shapes are identical (even if it is along the perpendicular bisector) (tasks \textit{1, 20}; see Figure~\ref{fig:exampleSD}), determining if all the shapes on an image are the same (\textit{16, 22}), or detecting if two pairs of identical shapes can be translated to become identical to each other
(\textit{13}). Another set of tasks within this sub-cluster includes tasks that are defined by more complex rules that involve the composition of additional relational judgments. Sample tasks include identifying pairs/triplets of identical shapes and measuring the distance with the rest (\textit{6, 17}), determining if an image consists of pairs of identical shapes (\textit{5}), or detecting if one of the shapes is a scaled version of the other (\textit{19}). Finally, the second sub-cluster $SD_1$ shown in dark-green color involves two tasks that require an understanding of shape transformations. One task asks observers to say if one of the shapes is the scaled, translated, or rotated version of the other one (\textit{21}). The other task test asks observers to judge whether or not an image contains two pairs of three identical shapes or three pairs of two identical shapes in an image (\textit{7}).  

To summarize this first set of experiments, we have systematically evaluated the ability of ResNets spanning multiple depths to solve each of the twenty-three SVRT tasks for different training set sizes. This allowed us to represent SVRT tasks according to their learnability by ResNets of varying depth. By clustering these representations, we extracted a novel SVRT taxonomy that both recapitulated an already described SD-SR dichotomy \citep{kim2018not}, and also revealed a more granular task structure corresponding to the number of rules used to form each task. Tasks with more rules are harder for ResNets to learn. Our taxonomy also reveals an organization of tasks where easier $SR_1$ and $SR_2$ sub-clusters fall closer to each other than harder $SD_1$ and $SD_2$ sub-clusters.

\section{Conclusion}

The goal of the present study was to shed light on the computational mechanisms underlying visual reasoning using the Synthetic Visual Reasoning Test (SVRT)~\citep{fleuret2011comparing}. \textcolor{black}{There are} twenty-three binary classification problems in this challenge, which include a variety of same-different and spatial reasoning tasks.  

In our experiment, we systematically evaluated the ability of a battery of $N=15$ deep convolutional neural networks (ResNets) -- varying in depths and trained using different training set sizes -- to solve each of the SVRT problems. We found a range of accuracies across all twenty-three tasks. Shallower networks easily learned some tasks, and relatively small training sets and some tasks were hardly solved with much deeper networks and orders of magnitude more training examples. 


Under the assumption that the computational complexity of individual tasks can be well characterized by the pattern of test accuracy across these $N=15$ neural networks, we formed N-dimensional accuracy vectors for each task and ran a hierarchical clustering algorithm. The resulting analysis \textcolor{black}{suggests} a taxonomy of visual reasoning tasks: beyond two primary clusters corresponding to same-different (SD) vs. spatial relation (SR) judgments, we also identified a finer organization with sub-clusters reflecting the nature and the number of relations used to compose the rules defining the task. Our results are consistent with previous work by \citet{kim2018not}, who first identified a dichotomy between SD and SR tasks. Our results also extend prior work \citep{fleuret2011comparing,kim2018not,yihe2019program} in \textcolor{black}{proposing} a finer-level taxonomy of visual reasoning tasks. The accuracy of neural networks is reflected in the number of relationships used to define the basic rules, which is expected, but it deserves closer examination. 

\citet{kim2018not} have previously suggested that SD tasks ``strain'' convolutional neural networks. That is, while it is possible to find a network architecture of sufficient depth (or the number of units) that can solve a version of the task up to a number of stimulus configurations (e.g., by forcing all stimuli to be contained within a $\Delta H \times \Delta W$ window), it is relatively easy to render the same task unlearnable by the same network \textcolor{black}{past} a certain number of stimulus configurations (e.g., by increasing the size of the window that contains all stimuli). It is as if these convolutional networks are \textcolor{black}{capable of learning} the task if the number of stimulus configurations remains below their memory capacity, and fails beyond that. \textcolor{black}{It remains an open question  whether non-convolution} alternatives to the CNNs tested here such as the now popular transformer networks  \citep{Dosovitskiy2020-iq,touvron2021training,tolstikhin2021mlp} would learn to solve some of the harder SVRT tasks more efficiently. As an initial experiment, we attempted to train and test a Vision Transformer \footnote[1]{\url{https://github.com/facebookresearch/dino}} (ViT) \citep{Dosovitskiy2020-iq}  constrained to have a similar number of parameters (21M) to the ResNet-50 used here. We were not able to get these architectures to do well on most of the tasks that are difficult for ResNets, even with 100k samples (also shown in \citet{messina2021recurrent}). It is worth noting that even 100k samples remain a relatively small dataset size by modern-day standards since ViT was trained from scratch.

Multi-layer perceptrons and convolutional neural networks including ResNets and other architectures can be formally shown to be universal approximators under certain architectural constraints. That is, they can learn arbitrary mappings between images to class labels. Depending on the complexity of the mapping, one might need an increasing number of hidden units to allow for enough expressiveness of the network; but provided enough units $/$ depth and a sufficient amount of  training examples, deep CNNs can learn arbitrary visual reasoning tasks. While we cannot make any strong claim for the specific ResNet architectures used in this study (currently, the proof is limited to a single layer without max pooling or batch normalization \citep{lin2018resnet}), we have indeed found empirically that all SVRT tasks could indeed be learned for networks of sufficient depth and provided a sufficient amount of training examples.
However, deep CNNs typically lack many of the human cognitive functions, such as attention and working memory. Such functions are likely to provide a critical advantage for a learner to solve some of these tasks \citep{10.7551/mitpress/1187.001.0001}. CNNs might have to rely instead on function approximation which could lead to a less general ``brute-force'' solution. 
Given this, an open question is whether the clustering of SVRT tasks derived from our CNN-based analyses will indeed hold for human studies. At the same time, the prediction by \citet{kim2018not} using CNNs that SD tasks are harder than SR tasks and hence that they may demand additional computations (through feedback processes) such as attention and/or working memory was successfully validated experimentally by \citet{AlamiaENEURO.0267-20.2020} using EEG. 

Additional evidence for the benefits of feedback mechanisms for visual reasoning was provided by  \citet{linsley2018learning} who showed that contour tracing tasks that can be solved efficiently with a single layer of a recurrent CNN may require several order of magnitudes more processing stages in a non-recurrent-CNN to solve the same task. This ultimately translates into much greater sample efficiency for recurrent-CNNs on natural image segmentation tasks \citep{linsley2020recurrent}. \textcolor{black}{The closely related task of ``insideness'' was also studied by \citet{villalobos2021neural} who demonstrated the inability of CNNs to learn a general solution for this class of problems.}
Universal approximators with minimal inductive biases such as multi-layer perceptrons, CNNs and other feedforward or \textcolor{black}{non}-attentive architectures can learn to solve visual reasoning tasks, but they might need a very large number of training examples to properly fit. Hence, beyond simply measuring the accuracy of very deep nets in high data regimes (such as when millions of training examples are available), systematically assessing the performance of neural nets of varying depths and for different training regimes may provide critical information about the complexity of different visual reasoning tasks.